% from ega1-10-fr.tex
\oldpage[I]{203}

Enfin, a) entraîne b). Considéré comme
$\mathscr{O}_{\mathfrak{X}}$-Module, on a
$\mathscr{O}_{X_n}=\mathscr{O}_{\mathfrak{X}}/
\mathscr{J}^{n+1}=A_n^\Delta$~; $\mathscr{F}_n=\mathscr{F}
\otimes_{\mathscr{O}_{\mathfrak{X}}}\mathscr{O}_{X_n}$ est un
$\mathscr{O}_{\mathfrak{X}}$-Module cohérent
(\hyperref[0.5.3.5-fr]{0, 5.3.5}), et comme c'est aussi un
$\mathscr{O}_{X_n}$-Module et que $\mathscr{J}^{n+1}$ est cohérent,
on en conclut que $\mathscr{F}_n$ est un
$\mathscr{O}_{X_n}$-Module cohérent
(\hyperref[0.5.3.10-fr]{0, 5.3.10}), et il est immédiat que
$u_{mn}^*(\mathscr{F}_n)=\mathscr{F}_m$ pour $m\leq n$ (en se
souvenant que l'application continue $X_m\to X_n$ des espaces
sous-jacents est l'identité de $\mathfrak{X}$). Le faisceau
$\mathscr{G}=\varprojlim\mathscr{F}_n$ est donc un
$\mathscr{O}_{\mathfrak{X}}$-Module cohérent, puisqu'on a vu que b)
entraîne a). Les homomorphismes canoniques
$\mathscr{F}\to\mathscr{F}_n$ forment un système projectif, qui par
passage à la limite donne un homomorphisme canonique
$w:\mathscr{F}\to\mathscr{G}$, et tout revient à démontrer que $w$ est
bijectif. La question étant maintenant locale, on peut se borner au cas
où $\mathscr{F}$ est conoyau d'un homomorphisme
$\mathscr{O}_{\mathfrak{X}}^p\to\mathscr{O}_{\mathfrak{X}}^q$~; cet
homomorphisme étant de la forme $v^\Delta$, où $v$ est un homomorphisme
$A^m\to A^n$ (\hyperref[I.10.10.2-fr]{10.10.2}), $\mathscr{F}$ est
isomorphe à $M^\Delta$, où $M=\operatorname{Coker}v$
(\hyperref[I.10.10.2-fr]{10.10.2}). On a alors, en vertu de
(\hyperref[I.10.10.2-fr]{10.10.2}),
$\mathscr{F}_n=M^\Delta\otimes_{\mathscr{O}_{\mathfrak{X}}}A_n^\Delta
=(M\otimes_A A_n)^\Delta$, et comme la topologie $\mathfrak{J}$-adique
sur $M\otimes_A A_n$ est discrète, on a
$(M\otimes_A A_n)^\Delta=(M\otimes_A A_n)^{\sim}$ (en tant que
$\mathscr{O}_{X_n}$-Module)~; on a vu plus haut que
$M^\Delta=\varprojlim\mathscr{F}_n$ et $w$ est donc bien dans ce cas
l'identité.
\par\hfill C.Q.F.D.\par

\begin{corollary}[10.10.6]
\label{I.10.10.6-fr}
Si $\mathscr{F}$ vérifie la condition b) de
(\hyperref[I.10.10.5-fr]{10.10.5}), le système projectif
$(\mathscr{F}_n)$ est isomorphe au système des
$\mathscr{F}\otimes_{\mathscr{O}_{\mathfrak{X}}}\mathscr{O}_{X_n}$.
\end{corollary}

\begin{env}[10.10.7]
\label{I.10.10.7-fr}
Soient maintenant $A$, $B$ deux anneaux adiques noethériens,
$\varphi:B\to A$ un homomorphisme continu~; on désignera par
$\mathfrak{J}$ (resp. $\mathfrak{K}$) un idéal de définition de $A$
(resp. $B$), tels que $\varphi(\mathfrak{K})\subset\mathfrak{J}$, et
on posera $X=\operatorname{Spec}(A)$, $Y=\operatorname{Spec}(B)$,
$\mathfrak{X}=\operatorname{Spf}(A)$,
$\mathfrak{Y}=\operatorname{Spf}(B)$. Soient $f:X\to Y$ le morphisme
de préschémas correspondant à $\varphi$
(\hyperref[I.1.6.1-fr]{1.6.1}), $\widehat f:\mathfrak{X}\to
\mathfrak{Y}$ son prolongement aux complétés
(\hyperref[I.10.9.1-fr]{10.9.1}), qui est aussi le morphisme de
préschémas formels correspondant à $\varphi$
(\hyperref[I.10.2.2-fr]{10.2.2}).
\end{env}

\begin{proposition}[10.10.8]
\label{I.10.10.8-fr}
Pour tout $B$-module $N$ de type fini, il existe un isomorphisme
canonique fonctoriel de $\mathscr{O}_{\mathfrak{X}}$-Modules
\[
  \widehat f^*(N^\Delta)
  \xrightarrow{\sim}
  (N\otimes_B A)^\Delta.
\]
\end{proposition}

En effet, en désignant par $i_X:\mathfrak{X}\to X$ et
$i_Y:\mathfrak{Y}\to Y$ les morphismes canoniques, on a
(\hyperref[I.10.8.8-fr]{10.8.8}), à des isomorphismes canoniques
fonctoriels près, $N^\Delta=i_Y^*(\widetilde N)$ et
\[
  (N\otimes_B A)^\Delta
  =i_X^*((N\otimes_B A)^{\sim})
  =i_X^*(f^*(\widetilde N))
\]
(\hyperref[I.1.6.5-fr]{1.6.5})~; la proposition résulte donc de la
commutativité du diagramme (\hyperref[I.10.9.2-fr]{10.9.2}).

\begin{corollary}[10.10.9]
\label{I.10.10.9-fr}
Pour tout idéal $\mathfrak b$ de $B$, on a
$\widehat f^*(\mathfrak b^\Delta)\mathscr{O}_{\mathfrak{X}}
=(\mathfrak bA)^\Delta$.
\end{corollary}

En effet, soit $j$ l'injection canonique $\mathfrak b\to B$, à laquelle
il correspond l'injection canonique
$j^\Delta:\mathfrak b^\Delta\to\mathscr{O}_{\mathfrak{Y}}$ de
faisceaux de $\mathscr{O}_{\mathfrak{Y}}$-Modules~; par définition,
$\widehat f^*(\mathfrak b^\Delta)\mathscr{O}_{\mathfrak{X}}$ est
l'image de l'homomorphisme
$\widehat f^*(j^\Delta):\widehat f^*(\mathfrak b^\Delta)\to
\mathscr{O}_{\mathfrak{X}}=\widehat f^*(\mathscr{O}_{\mathfrak{Y}})$~;
mais cet homomorphisme s'identifie à
$(j\otimes1)^\Delta:(\mathfrak b\otimes_B A)^\Delta\to
\mathscr{O}_{\mathfrak{X}}=(B\otimes_B A)^\Delta$ d'après
(\hyperref[I.10.10.8-fr]{10.10.8}). Comme l'image de $j\otimes1$ est
l'idéal $\mathfrak bA$ de $A$, l'image de $(j\otimes1)^\Delta$ est
donc $(\mathfrak bA)^\Delta$ en vertu de
(\hyperref[I.10.10.2-fr]{10.10.2}), d'où la conclusion.
